% !TeX root = ../cosmology_summary.tex

%% --------------------------------------------------------
\chapter{Ефект Закса-Вольфа та формування кутового спектра CMB}
\label{sec:sachs_wolfe_main}
%% --------------------------------------------------------

Кутова анізотропія температури реліктового випромінювання (CMB), яку ми спостерігаємо сьогодні (рис.~\ref{fig:planck_map}), є безпосереднім реліктовим \enquote{знімком} стану раннього Всесвіту в епоху рекомбінації ($z_{\mathrm{rec}} \approx 1100$). Флуктуації температури на небесній сфері формуються внаслідок накладання кількох фізичних процесів, що відбувалися на поверхні останнього розсіяння (LSS).


\begin{figure}[htbp]
\centering
\includegraphics[width=\textwidth]{planck_cmb_map.png}
\caption{Карта температурних анізотропій CMB від місії
         Planck~2018 \cite{planck_2018}. Кольори відображають
         відхилення від середнього $T_0 = 2.725\,\mathrm{K}$
         у діапазоні $\pm 300\,\mu\mathrm{K}$~--- спостережуваний
         відбиток первинних флуктуацій інфлатона після
         ${\sim}380\,000$ років еволюції.
         Зображення: ESA and the Planck Collaboration
         \cite{esa_planck_cmb_image}.}
\label{fig:planck_map}
\end{figure}

З погляду релятивістської кінетичної теорії, повна варіація температури у заданому напрямку $\hat{\mathbf{n}}$ описується калібровочно-інваріантним рівнянням, що враховує як внутрішні властивості плазми, так і геометрію простору-часу \cite{hu_sugiyama_1995}:
\begin{equation}\label{eq:cmb_sources_combined}
  \frac{\delta T}{T}(\hat{\mathbf{n}}) = \left( \Theta_0 + \Phi \right)_{\mathrm{rec}} + \left( \hat{\mathbf{n}} \cdot \mathbf{v}_b \right)_{\mathrm{rec}} + 2 \int_{\eta_{\mathrm{rec}}}^{\eta_0} \frac{\partial \Phi}{\partial \eta} \, d\eta,
\end{equation}
де $\Theta_0 = \delta_\gamma/4 = \delta T/T$ на поверхні останнього розсіяння (див.~\eqref{eq:Theta0_def}).

Тут перший доданок відображає \emph{ефект Закса-Вольфа}~---
комбінацію монопольної компоненти збурення температури фотонів
$\Theta_0$ (введеної у розділі~\ref{sec:bao})
та гравітаційного червоного зміщення $\Phi$ при виході фотонів
з потенціальних ям. Другий доданок відповідає за доплерівський
зсув частоти через пекулярний рух баріонної компоненти зі
швидкістю $\mathbf{v}_b$. Третій доданок є \emph{інтегральним
ефектом Закса-Вольфа} (ISW)~--- він враховує еволюцію
гравітаційних потенціалів $\dot{\Phi}$ вздовж траєкторії
фотонів до спостерігача. ISW діє на двох епохах: у ранній
радіаційно-домінантній фазі потенціали розпадаються через
перехід мод під горизонт, де тиск радіації підтримує
коливання, що дає ранній ISW ефект на мультиполях
$\ell \sim 100$--$200$; а в сучасну епоху темна енергія
прискорює розширення і потенціали знову змінюються~---
це пізній ISW, який проявляється як слабкий підйом спектра
при найменших $\ell \lesssim 10$.

Рівняння \eqref{eq:cmb_sources_combined} є фундаментальним, оскільки воно повністю описує локальну мікрофізику на реліктовому обрії. Проте безпосереднє зіставлення цього виразу з астрофізичними спостереженнями стикається з принциповою перешкодою: ми спроможні виміряти параметри плазми не в кожній тривимірній точці раннього Всесвіту, а лише фіксуємо двовимірну проекцію цих процесів на небесну сферу. Щобільше, конкретний візерунок гарячих і холодних плям у нашому куточку Космосу є випадковою реалізацією первинних квантових флуктуацій. Для перевірки теоретичних моделей нам потрібні не локальні значення $\delta T / T$ у конкретному напрямку, а об'єктивні статистичні характеристики, що є інваріантними відносно положення спостерігача. Саме ця потреба змушує перейти від конфігураційного простору до гармонічного аналізу сигналів на сфері.

%% --------------------------------------------------------
\section{Гармонічний аналіз та кутовий спектр потужності}
\label{subsec:harmonic_analysis}
%% --------------------------------------------------------

Оскільки спостережуваний візерунок температурних флуктуацій $\Delta T(\theta,\phi) = T(\theta,\phi) - T_\text{avg} $ фіксується на двовимірній сфері, природним математичним базисом для його статистичного опису є розклад за ортонормованими \emph{сферичними гармоніками} $Y_{\ell m}(\theta,\phi)$:
\begin{equation}\label{eq:spherical_decomp_integrated}
    \frac{\Delta T}{T}(\theta,\phi) = \sum_{\ell=0}^{\infty}\sum_{m=-\ell}^{\ell} a_{\ell m}\,Y_{\ell m}(\theta,\phi),
\end{equation}
де безрозмірні комплекснозначні коефіцієнти $a_{\ell m}$ відіграють роль амплітуд флуктуацій. Мультипольний момент $\ell$ визначає просторову частоту і є обернено пропорційним кутовому масштабу структури на небі: $\theta \approx 180^\circ/\ell$.

Для побудови теорії, вільної від вибору конкретної системи координат чи орієнтації інструментів спостереження, вводять безрозмірний \emph{кутовий спектр потужності} $C_\ell$. Припускаючи статистичну ізотропію та однорідність первинних збурень, двоточковий корелятор коефіцієнтів усереднюється за всіма можливими проекціями $m$, повністю затираючи просторові координати:
\begin{equation}
    \langle a_{\ell m} a_{\ell' m'}^* \rangle = C_\ell \, \delta_{\ell \ell'} \, \delta_{m m'}.
\end{equation}

Проте в реальних експериментах (наприклад, у космічній місії \emph{Planck}) вимірюють флуктуації безпосередньо в термодинамічних одиницях температури (мікрокельвінах, $\mu\text{K}$). Крім того, на практиці замість сирого спектра $C_\ell$ прийнято оперувати \emph{масштабованим (або згладженим) кутовим спектром потужності} $\mathcal{D}_\ell^{TT}$, який означується як:
\begin{equation}\label{eq:Dl_definition}
    \mathcal{D}_\ell^{TT} \equiv \frac{\ell(\ell+1)}{2\pi} T_0^2 \, C_\ell^{TT},
\end{equation}
де $T_0 \approx 2.7255\text{ К}$ --- середня температура реліктового фону сьогодні. Множник $\ell(\ell+1)$ має глибокий фізичний зміст: він повністю компенсує спадання амплітуди великомасштабних флуктуацій, перетворюючи чисто геометричний ефект Закса-Вольфа на плоске плато в області малих мультиполів, що значно спрощує візуальний аналіз спектра.

Прямий фізичний зв'язок між тривимірними модами первинних збурень у просторі Фур'є (що мають хвильове число $k$) та двовимірними мультиполями $\ell$ задається відстанню кутового діаметра $d_A$ до поверхні останнього розсіяння. З простої геометричної проекції плоскої хвилі довжиною $\lambda = 2\pi/k$ випливає наближене співвідношення:
\begin{equation}\label{eq:l_k_projection_main}
    \ell \approx k\cdot d_A.
\end{equation}

Строгий математичний перехід здійснюється через інтегрування за допомогою сферичних функцій Бесселя $j_\ell(kd_A)$:
\begin{equation}\label{eq:bessel_integral_main}
    C_\ell \propto \int_0^{\infty}|\Theta_k(\eta_\text{rec})|^2\cdot j_\ell^2(kd_A)\cdot k^2\,dk.
\end{equation}
Кінематична властивість функції $j_\ell(x)$ полягає в тому, що вона є практично нульовою при $x < \ell$ і зазнає різкого локального сплеску саме в околі точки $x \approx \ell$, що математично легалізує та обґрунтовує наближення \eqref{eq:l_k_projection_main}.

%% --------------------------------------------------------
\section{Від акустичних коливань до піків спектра}
\label{subsec:acoustic_oscillations}
%% --------------------------------------------------------

У розділі~\ref{sec:bao} ми показали, що фотон-баріонна
плазма на субхабблових масштабах коливається як
гармонічний осцилятор~\eqref{eq:Theta0_osc}, і що
розв'язок для адіабатичних початкових умов має вигляд
$\Theta_0(k,\eta) \approx \Theta_0(0)\cos(kc_s\eta)$
(формула~\eqref{eq:Theta0_solution}).

Рекомбінація діє як «стоп-кадр»: фаза коливань
заморожується при $\eta = \eta_{\rm rec}$.
Внесок кожної моди $k$ у кутову анізотропію визначається
саме цією замороженою амплітудою, тому кутовий спектр
потужності пропорційний її квадрату:
\begin{equation}\label{eq:Cl_cos2}
  C_\ell \propto \cos^2\!\left(kc_s\eta_{\rm rec}\right)
  \Big|_{k \approx \ell/d_A}.
\end{equation}
Екстремуми при $kc_s\eta_{\rm rec} = \pi, 2\pi, 3\pi,\ldots$
дають серію \emph{акустичних піків} у спектрі
$\mathcal{D}_\ell^{TT}$.

На супергоризонтних масштабах ($k \ll aH$, тобто
$\ell \lesssim 30$) коливань немає~--- домінує чисто
геометричний ефект Закса-Вольфа.
Для адіабатичних збурень $\Theta_0 = -\tfrac{2}{3}\Phi$,
і ефективна температурна флуктуація при виході з ями:
\begin{equation}
  \left(\Theta_0 + \Phi\right)\Big|_{\text{супергоризонт}}
  = \frac{1}{3}\Phi = \mathrm{const},
\end{equation}
що після множення на $\ell(\ell+1)$ у
формулі~\eqref{eq:Dl_definition} дає \emph{горизонтальне
плато} при малих $\ell$.

Важливо: звукові хвилі не затерли первинний розподіл
густини від інфлатона, а лише \emph{модулювали} його.
Темна матерія не взаємодіє з фотонами~--- її потенціальні
ями залишалися нерухомими, і баріон-фотонний газ
коливався відносно цих застиглих центрів.
Тому $\Theta_0(0)$ несе в собі первинний квантовий
«відбиток» інфлатона.

%% --------------------------------------------------------
\section{Ефект баріонного завантаження та порушення симетрії}
\label{subsec:baryon_loading_effects}
%% --------------------------------------------------------

Наявність важкої баріонної компоненти, що описується параметром $R$, відображається на внутрішній симетрії фотонних осциляцій. Якщо повернутися до точного розв'язку неоднорідного рівняння \eqref{eq:Theta0_osc} зі збереженням гравітаційного члена у правій частині, виявляється, що присутність баріонів зміщує положення динамічної рівноваги (нуля) коливань осцилятора. Математично, загальний розв'язок для монополя температури набуває вигляду:
\begin{equation}\label{eq:shifted_oscillator_final_fixed}
    \Theta_0(\eta_\text{rec}) = \Theta_0(0)\cdot\cos\!\left(kc_s\eta_\text{rec}\right) - (1+R)\Phi,
\end{equation}
де $\Theta_0(0)$ --- амплітуда косинусних коливань (початкове значення у $\eta=0$), а $(1+R)\Phi$ --- зміщення рівноваги, зумовлене гравітаційним потенціалом.

З урахуванням ефекту Закса-Вольфа, повне ефективне збурення, яке виривається з потенціальної ями, визначається комбінацією $\Theta_0 + \Phi$. Підставляючи сюди розв'язок \eqref{eq:shifted_oscillator_final_fixed}, знаходимо:
\begin{equation}\label{eq:baryon_loading_shift_result}
    \left( \Theta_0 + \Phi \right)_{\mathrm{rec}} = \Theta_0(0)\cdot\cos\!\left(kc_s\eta_\text{rec}\right) - R\Phi.
\end{equation}

Безрозмірний зсув нуля $\Delta \equiv -R\Phi / \Theta_0(0)$
створює суттєву асиметрію між парними та непарними акустичними
піками у спектрі потужності $\mathcal{D}_\ell^{TT}$.
Оскільки $\Phi < 0$ у потенціальній ямі темної матерії
(гравітаційний потенціал від'ємний за визначенням), а $R > 0$,
маємо $\Delta > 0$~--- зсув рівноваги осцилятора відбувається
у бік стиснення:
\begin{itemize}
    \item \textbf{Непарні піки (перший, третій тощо)} відповідають фазі максимального гравітаційного стиснення плазми всередині потенціальних ям темної матерії. Гравітаційне поле баріонів діє синергетично з колапсом, підсилюючи амплітуду ефективного збурення до величини $(1+\Delta)$. Відповідно, кутова потужність пропорційна квадрату цього значення:
    \begin{equation}
        \text{Потужність непарних піків} \propto (1+\Delta)^2.
    \end{equation}

    \item \textbf{Парні піки (другий, четвертий тощо)} фіксують протилежну фазу --- момент максимального пружного розширення (радіаційного розрідження). У цій точці фотонний тиск намагається викинути речовину з ями, але гравітаційний потенціал баріонів утримує її, послаблюючи чисту амплітуду до $(-1+\Delta)$. Для кутового спектра потужності це дає:
    \begin{equation}
        \text{Потужність парних піків} \propto (-1+\Delta)^2 = (1-\Delta)^2.
    \end{equation}
\end{itemize}

Оскільки і метричний потенціал $\Phi$, і початкове збурення плазми $\Theta_0(0)$ народилися з єдиного первинного скалярного поля, їхнє відношення є безрозмірною константою порядку одиниці. Якби у Всесвіті були відсутні баріони ($R=0 \implies \Delta = 0$), осцилятор коливався б симетрично відносно нуля, і перший та другий піки мали б абсолютно однакову висоту. Отже, прецизійне вимірювання відносної висоти перших двох піків спектра $\mathcal{D}_\ell^{TT}$ дозволяє астрофізикам з високою точністю \enquote{зважити} баріонний вміст Всесвіту, відокремивши його від внеску холодної темної матерії (див. рис.~\ref{fig:cmb_tt}).

\begin{figure}[htbp]
\centering
\begin{tikzpicture}
\begin{semilogxaxis} [
    width    = \linewidth,
    height  = \linewidth,
    xlabel  = {Мультипольний момент $\ell$},
    ylabel  = {$\mathcal{D}_\ell^{TT}\ [\mu\mathrm{K}^2]$},
    yticklabel style={/pgf/number format/1000 sep={\,}},
    xmin    = 2,   xmax = 2600,
    ymin    = 0,   ymax = 6800,
    % --- РЯТУЄМО ОСЬ Х ПІД \linewidth ---
    xtick = {2, 10, 30, 100, 500, 1000, 1500, 2000, 2500},
    % Використовуємо компактний науковий формат, щоб мітки не зливалися на вузькому графіку:
    xticklabels = {2, 10, 30, 100, 500, , 1500},
    % Трохи зменшуємо шрифт міток для осі Х та зміщуємо вниз, щоб сітка не заважала
    ticklabel style = {font=\footnotesize, yshift=-1pt},
    % Густа логарифмічна сітка
    minor xtick = {3,4,5,6,7,8,9,20,40,50,60,70,80,90,200,300,400,600,700,800,900,1100,1200,1300,1400,1600,1700,1800,1900,2100,2200,2300,2400},
    ytick   = {0, 1000, 2000, 3000, 4000, 5000, 6000},
    grid   = both,
    grid style = {line width=0.3pt, draw=gray!25},
    major grid style = {line width=0.5pt, draw=gray!40},
    label style      = {font=\small},
    legend style = {
        at={(0.6,0.97)}, anchor=north east,
        font=\small, fill=white, fill opacity=0.9,
        draw=gray!60, inner sep=4pt, row sep=1pt
    },
    clip = true,
]

    %--- LambdaCDM Best-fit Curve ---
    \addplot[
        color     = theoryred,
        width     = 1.6pt,
        smooth,
    ] table[x=ell, y=Dl] {\theorydata};
    \addlegendentry{$\Lambda$CDM (CAMB, Planck\,2018 best-fit)}

    %--- Planck 2018 data with error bars ---
    \addplot[
        only marks,
        mark       = *,
        mark size  = 1.4pt,
        color      = planckblue,
        error bars/.cd,
            y dir  = both,
            y explicit,
    ] table[
        x    = ell,
        y    = Dl,
        y error minus = errm,
        y error plus  = errp
    ] {\planckdata};
    \addlegendentry{Planck\,2018 (збіновані дані TT)}

    %--- Annotate acoustic peaks ---
    \node[font=\tiny, text=gray!70, above] at (axis cs:220,5732) {1-й пік};
    \node[font=\tiny, text=gray!70, above] at (axis cs:540,2592) {2-й пік};
    \node[font=\tiny, text=gray!70, above] at (axis cs:810,2540) {3-й пік};

    %--- Sachs-Wolfe plateau annotation ---
    \draw[->, gray!60, thin] (axis cs:6,1300) -- (axis cs:30,1100);
    \node[font=\tiny, text=gray!70, anchor=south, align=center] at (axis cs:12,1350) {Плато\\Закса-Вольфа};

  \end{semilogxaxis}
\end{tikzpicture}
\caption{Кутовий спектр потужності температурної анізотропії реліктового випромінювання за даними космічної місії \emph{Planck 2018}. Суцільна червона лінія відповідає найкращему фіту стандартної космологічної моделі $\Lambda$CDM. Положення першого акустичного піка ($\ell \approx 220$) жорстко фіксує кутовий розмір звукового горизонту на момент рекомбінації, виступаючи чутливим індикатором просторової кривини Космосу ($\Omega_k \approx 0$). Відносна висота парних і непарних піків відображає ефект баріонного завантаження: відношення другого піка до першого дозволяє з високою точністю обчислити фізичну густину баріонів $\Omega_b h^2$, тоді як третій пік контролює густину холодної темної матерії $\Omega_c h^2$.}
\label{fig:cmb_tt}
\end{figure}

\begin{figure}[htbp]
\centering
\begin{tikzpicture}
\begin{axis}[
    width   = \linewidth,
    height  = \linewidth,
    xlabel  = {Мультипольний момент $\ell$},
    ylabel  = {$\mathcal{D}_\ell^{TT}\ [\mu\mathrm{K}^2]$},
    xmin    = 2,   xmax = 80,
    ymin    = 650, ymax = 1150,
    xtick   = {2, 10, 20, 30, 40, 50, 60, 70, 80},
    ytick   = {700, 800, 900, 1000, 1100},
    grid    = both,
    grid style       = {line width=0.3pt, draw=gray!25},
    major grid style = {line width=0.5pt, draw=gray!40},
    label style      = {font=\small},
    ticklabel style  = {font=\footnotesize},
    legend style = {
        at={(0.98,0.97)}, anchor=north east,
        font=\small, fill=white, fill opacity=0.9,
        draw=gray!60, inner sep=4pt, row sep=1pt
    },
    clip = true,
]

    %--- LambdaCDM Best-fit Curve ---
    \addplot[
        color = theoryred,
        line width = 1.6pt,
        smooth,
    ] table[x=ell, y=Dl] {\theorydata};
    \addlegendentry{$\Lambda$CDM (CAMB)}

    %--- Planck 2018 data with error bars ---
    \addplot[
        only marks,
        mark      = *,
        mark size = 1.8pt,
        color     = planckblue,
        error bars/.cd,
            y dir = both,
            y explicit,
    ] table[
        x              = ell,
        y              = Dl,
        y error minus  = errm,
        y error plus   = errp,
    ] {\planckdata};
    \addlegendentry{Planck\,2018}

    %--- n_s = 1 horizontal reference (Harrison-Zeldovich) ---
    % Normalised at ell=10, D_ell=819
    \addplot[
        domain  = 2:80,
        samples = 200,
        dashed,
        line width = 1.2pt,
        color   = gray!60,
    ] {819};
    \addlegendentry{$n_s = 1$ (масштабоінваріантний)}

    %--- n_s = 0.9649 (Planck best-fit) ---
    \addplot[
        domain  = 2:80,
        samples = 200,
        dashed,
        line width = 1.5pt,
        color   = orange!90!black,
    ] {819 * (x/10)^(0.9649 - 1)};
    \addlegendentry{$n_s = 0.965$ (Planck\,2018)}

    %--- n_s = 0.94 ---
    \addplot[
        domain  = 2:80,
        samples = 200,
        dashed,
        line width = 1.2pt,
        color   = teal,
    ] {819 * (x/10)^(0.94 - 1)};
    \addlegendentry{$n_s = 0.94$}

    %--- Annotation: formula ---
    \node[font=\small, text=orange!80!black, anchor=south west]
        at (axis cs:3, 660)
        {$\mathcal{D}_\ell^{TT} \propto \ell^{\,n_s - 1}$};

    %--- Arrow pointing to SW plateau ---
    \draw[->, gray!60, thin] (axis cs:6, 1100) -- (axis cs:25, 1010);
    \node[font=\tiny, text=gray!70, anchor=south, align=center]
        at (axis cs:10, 1110)
        {Плато\\Закса-Вольфа};

\end{axis}
\end{tikzpicture}
\caption{%
    Збільшений вигляд плато Закса--Вольфа ($2 \leq \ell \lesssim 80$)
    кутового спектра потужності $\mathcal{D}_\ell^{TT}$
    за даними місії \emph{Planck\,2018}~\cite{planck_2018}.
    Штрихові лінії ілюструють нахил первинного спектра
    $\mathcal{D}_\ell^{TT} \propto \ell^{\,n_s - 1}$
    для трьох значень спектрального індексу~$n_s$.
    Горизонтальна лінія відповідає масштабоінваріантному спектру
    Гарісона--Зельдовича ($n_s = 1$);
    відхилення від неї для планківського значення
    $n_s = 0.965 \pm 0.004$ підтверджує \emph{червоний нахил}
    первинних скалярних збурень~--- базове передбачення
    інфляційних моделей повільного кочення поля.
    Усі криві нормовано в точці $\ell = 10$,
    де $\mathcal{D}_{10}^{\Lambda\mathrm{CDM}} \approx 819\,\mu\mathrm{K}^2$.%
}
\label{fig:ns_tilt}
\end{figure}

%% --------------------------------------------------------
\section*{Задачі}
\addcontentsline{toc}{section}{Задачі}
%% --------------------------------------------------------

\begin{problem}{Ефект Закса-Вольфа та супергоризонтні масштаби}
Для адіабатичних збурень на супергоризонтних масштабах
($k \ll aH$) справедливе співвідношення $\Theta_0 = -\frac{2}{3}\Phi$.

\begin{enumerate}
    \item Покажіть, що комбінація $(\Theta_0 + \Phi)$ на цих масштабах
          дорівнює $\frac{1}{3}\Phi$ і не залежить від часу.
    \item Поясніть, чому цей результат призводить до горизонтального
          плато в масштабованому кутовому спектрі $\mathcal{D}_\ell^{TT}$
          при малих $\ell$ ($\ell \lesssim 30$).
    \item Як зміниться цей висновок, якщо збурення не є адіабатичними?
\end{enumerate}
\end{problem}

% -------------------------------------------------------------------

\begin{problem}{Від хвильового числа $k$ до мультиполя $\ell$}
Кутовий розмір звукового горизонту на момент рекомбінації
задається співвідношенням $\theta_s = r_s / d_A$, де
$r_s \approx 147$ Мпк --- звуковий горизонт, а
$d_A$ --- кутова відстань до поверхні останнього розсіяння.

\begin{enumerate}
    \item У плоскому Всесвіті ($\Omega_k = 0$) кутова відстань
          дорівнює $d_A = \eta_0 - \eta_{\text{rec}} \approx 14$ Гпк.
          Оцініть кутовий розмір звукового горизонту $\theta_s$ в радіанах
          та градусах.
    \item Використовуючи наближення $\ell \approx \pi / \theta_s$,
          оцініть мультипольний номер $\ell$, що відповідає першому
          акустичному піку. Порівняйте зі спостережуваним значенням
          $\ell \approx 200$.
    \item Чому наближення $\ell = k d_A$ працює добре для великих $\ell$,
          але дає похибку для малих?
\end{enumerate}
\end{problem}

% -------------------------------------------------------------------

\begin{problem}{Баріонне завантаження та асиметрія піків}
Точний розв'язок рівняння для $\Theta_0$ з урахуванням гравітаційного
потенціалу дає:
\[
(\Theta_0 + \Phi)_{\text{rec}} = \Theta_0(0)\cos(k r_s) - R\Phi,
\]
де $r_s = \int_0^{\eta_{\text{rec}}} c_s\,d\eta$ --- звуковий горизонт,
$R = 3\rho_b/4\rho_\gamma$ --- баріонне завантаження.

\begin{enumerate}
    \item Для парних піків $\cos(k r_s) = +1$, для непарних $\cos(k r_s) = -1$.
          Покажіть, що відношення потужностей непарного та парного піків дорівнює
          \[
          \frac{\mathcal{D}_{\text{непар}}}{\mathcal{D}_{\text{пар}}} =
          \left(\frac{1+\Delta}{1-\Delta}\right)^2,
          \]
          де $\Delta = -R\Phi/\Theta_0(0) > 0$.
    \item Використовуючи $\Phi \approx -2/3\,\mathcal{R}$ та
          $\Theta_0(0) \approx \mathcal{R}/3$ (для адіабатичних збурень),
          оцініть $\Delta$ через параметр $R$.
    \item Для $R \approx 0.6$ (значення на рекомбінації) обчисліть
          відношення висот першого (непарного) та другого (парного) піків.
          Порівняйте з рис.~\ref{fig:cmb_tt}.
\end{enumerate}
\end{problem}

% -------------------------------------------------------------------

\begin{problem}{Спектральний індекс та нахил плато}
У наближенні Закса-Вольфа ($\ell \lesssim 30$) масштабований
кутовий спектр $\mathcal{D}_\ell^{TT}$ пов'язаний із первинним
спектром $\mathcal{P}_\zeta(k)$ співвідношенням:
\[
\mathcal{D}_\ell^{TT} \propto \ell(\ell+1) C_\ell^{SW} \propto \mathcal{P}_\zeta(k)\Big|_{k \approx \ell/d_A}.
\]

\begin{enumerate}
    \item Припускаючи степеневий вигляд $\mathcal{P}_\zeta(k) = A_s (k/k_0)^{n_s-1}$,
          покажіть, що $\mathcal{D}_\ell^{TT} \propto \ell^{\,n_s-1}$.
    \item Використовуючи планківське значення $n_s = 0.965 \pm 0.004$,
          обчисліть відношення $\mathcal{D}_{30}^{TT} / \mathcal{D}_{10}^{TT}$.
    \item Чому відхилення $n_s$ від одиниці (так званий «червоний нахил»)
          є важливим підтвердженням інфляційної парадигми?
\end{enumerate}
\end{problem}

% -------------------------------------------------------------------

\begin{problem}{Інтегральний ефект Закса-Вольфа (ISW)}
Пізній ISW виникає, коли гравітаційний потенціал $\Phi$ змінюється
в часі через домінування темної енергії:
\[
\frac{\Delta T}{T}\Big|_{\text{ISW}} = 2 \int_{\eta_{\text{rec}}}^{\eta_0} \dot{\Phi}\,d\eta.
\]

\begin{enumerate}
    \item У матеріально-домінованому Всесвіті без темної енергії
          $\Phi = \text{const}$. Чому тоді ISW-внесок відсутній?
    \item Поясніть, чому пізній ISW проявляється лише на найменших
          мультиполях ($\ell \lesssim 10$).
    \item Запропонуйте спостережний тест, який дозволяє відокремити
          ISW-внесок від інших ефектів (наприклад, за допомогою
          крос-кореляції з даними великомасштабної структури).
\end{enumerate}
\end{problem}